\begin{figure*}[ht!]
	\centering
		\small
	\begin{PromptBox}[System Message]
		You are an expert systematic review information specialist.
		
		You are tasked to formulate a systematic review Boolean query step by step as a reasoning process within <think> </think>, and provide the Boolean query formulated <answer> </answer>.
	\end{PromptBox}
	
	\vspace{1em}
	
	\begin{PromptBox}[User Message]
		You are given a systematic review research topic, with the topic title "{topic}".
		
		You need to simulate a Boolean query construction process using the \textbf{objective method}, which is grounded in domain expertise and structured logic.
		
		\textbf{Step 1}: Simulate a concise title and abstract (2–3 sentences) of a \textit{relevant and focused} article clearly aligned with the topic. This is a hypothetical but plausible example.
		
		\textbf{Step 2}: Based on the simulated text, identify \textit{key informative terms or phrases} that best represent the article's core concepts. Prioritise specificity and informativeness. Avoid overly broad or ambiguous terms.
		
		\textbf{Step 3}: Categorise each term into one of the following:
		- (A) Health conditions or populations (e.g., diabetes, adolescents)
		- (B) Treatments, interventions, or exposures (e.g., insulin therapy, air pollution)
		- (C) Study designs or methodologies (e.g., randomized controlled trial, cohort study)
		- (N/A) Not applicable to any of the above categories
		
		\textbf{Step 4}: Using the categorised terms, build a Boolean query in MEDLINE format for PubMed:
		- Combine synonyms or related terms within each category using OR
		- Use both free-text terms and MeSH terms (e.g., chronic pain[tiab], Pain[mh])
		- \textbf{Do not wrap terms or phrases in double quotes}, as this disables automatic term mapping (ATM)
		- Tag each term individually when needed (e.g., covid-19[ti] vaccine[ti] children[ti])
		- Field tags limit the search to specific fields and disable ATM
		
		\textbf{Step 5}: Use wildcards (*) to capture word variants (e.g., vaccin* -> vaccine, vaccination):
		- Terms must have >= 4 characters before the * (e.g., colo*)
		- Wildcards work with field tags (e.g., breastfeed*[tiab]).
		
		\textbf{Step 6}: Combine all category blocks using AND:
		((itemA1[tiab] OR itemA2[tiab] OR itemA3[mh]) AND (itemB1[tiab] OR ...) AND (itemC1[tiab] OR ...))
		
		\textbf{Only use the following allowed field tags:}
		Title: [ti], Abstract: [ab], Title/Abstract: [tiab]
		MeSH: [mh], Major MeSH: [majr], Supplementary Concept: [nm]
		Text Words: [tw], All Fields: [all]
		Publication Type: [pt], Language: [la]
		
		Place your full reasoning (including simulated abstract, term list, classification, and query construction) inside <think></think>.
		
		Output the final Boolean query inside <answer></answer>.
		
		Do not include anything outside the <think> and <answer> tags.
		
		Do not include date restrictions.
	\end{PromptBox}
	
	\caption{Objective Reasoning prompt with simulated article and structured reasoning in <think>}
	\label{fig:prompt-objective}
\end{figure*}